\section{Conclusion}

FraudShiftBench treats an empirical result as a claim under a deployment contract, not as a score attached only to a model and dataset. This framing changes the scientific conclusion in the validated evidence. On Elliptic, a graph-visibility intervention changes the leading AUPRC model from GraphSAGE to GCN; on DGraphFin, GraphSAGE remains first but loses $21.5\%$ relative AUPRC. IBM AML-Data shows a second kind of instability: ranking and fixed-threshold metrics select different configurations in 11 of 16 feasible contexts, and construction tradeoffs change with scale. Resource guards leave Large and Medium GINE outside predictive ordering rather than assigning them implicit losses.

The support validator makes these boundaries executable. It rejected incomplete, prediction-missing, integrity-failed, construct-invalid, resource-widened, and contradicted claims in all 14 controlled cases. It also exposed three metadata labels that represented duplicate runs rather than distinct temporal stresses. The GraphSafe case illustrates the same discipline at the decision level: DGraphFin summaries are favorable, Elliptic favors simple averaging, and no comparison survives Holm correction across the declared 48-test family.

The central result is therefore conditional but actionable. Fraud benchmarks should report the contract that produced a ranking, separate ranking from operating-point utility, and keep unmeasured cells visible without turning them into performance evidence.
